\documentclass[a4paper,9pt]{extarticle}

%=========================================================
% PACKAGES
%=========================================================

\usepackage[
a4paper,
top=1.8cm,
bottom=1.5cm,
footskip=7mm,
left=1cm,
right=1cm,
headheight=14pt,
headsep=5mm
]{geometry}

\usepackage[T1]{fontenc}
\usepackage[utf8]{inputenc}
\usepackage[french]{babel}

\usepackage[sfdefault]{FiraSans}
\renewcommand{\familydefault}{\sfdefault}

\usepackage{multicol}
\setlength{\columnsep}{0.7cm}

\usepackage{amsmath,amssymb}
\usepackage{titlesec}
\usepackage{enumitem}

\usepackage{tcolorbox}
\tcbuselibrary{skins,breakable}

\usepackage{xcolor}
\usepackage{graphicx}
\usepackage{float}

\usepackage{fancyhdr}
\usepackage[hidelinks]{hyperref}


%=========================================================
% PAGE SETTINGS
%=========================================================

\setlength{\parindent}{0pt}
\setlength{\parskip}{1.5pt}

\setlength{\abovedisplayskip}{4pt}
\setlength{\belowdisplayskip}{4pt}

\setlist[itemize]{
leftmargin=*,
itemsep=1pt,
topsep=1pt
}
\setlist[enumerate]{
leftmargin=*,
itemsep=1pt,
topsep=1pt
}


%=========================================================
% COLORS
%=========================================================

\definecolor{Blue}{HTML}{E3F0FF}
\definecolor{Green}{HTML}{E3F7EA}
\definecolor{Pink}{HTML}{FFE5ED}
\definecolor{Orange}{HTML}{FFF0DC}
\definecolor{Purple}{HTML}{EEE7FF}
\definecolor{Yellow}{HTML}{FFFADB}
\definecolor{Goat}{HTML}{F7EEDF}

\definecolor{BlueTitle}{HTML}{285B8F}
\definecolor{GreenTitle}{HTML}{287052}
\definecolor{PinkTitle}{HTML}{A84063}
\definecolor{OrangeTitle}{HTML}{A96500}
\definecolor{PurpleTitle}{HTML}{6245A5}
\definecolor{GoatTitle}{HTML}{8A5A1F}
\definecolor{RecapTitle}{HTML}{5B3A8E}


%=========================================================
% HEADER
%=========================================================

\pagestyle{fancy}
\fancyhf{}

\fancyhead[L]{\small\textbf{Comportement mécanique des matériaux}}
\fancyhead[C]{\small Céramiques : panorama et structure}
\fancyhead[R]{\small Week 04 --- 28/09-04/10/2026}

\fancyfoot[C]{\small\href{https://github.com/AlphaXT38G/notes_ba5}{github.com/AlphaXT38G/notes\_ba5}}
\fancyfoot[R]{\small \thepage}

\renewcommand{\headrulewidth}{0.4pt}


%=========================================================
% TITLES
%=========================================================

\titleformat{\section}
{\large\bfseries\color{BlueTitle}}
{}{0pt}{}

\titlespacing{\section}
{0pt}{6pt}{3pt}


%=========================================================
% BOXES
%=========================================================
% Boxes are kept whole (unbreakable): inside multicols, a box split by tcolorbox
% leaves an empty column behind it.

\newtcolorbox{definition}{
enhanced, unbreakable, colback=Blue, colframe=BlueTitle,
boxrule=0.7pt, arc=0mm, title=\textbf{Définition}, fonttitle=\bfseries
}

\newtcolorbox{important}{
enhanced, unbreakable, colback=Pink, colframe=PinkTitle,
boxrule=0.7pt, arc=0mm, title=\textbf{Idée clé}, fonttitle=\bfseries
}

\newtcolorbox{formula}{
enhanced, unbreakable, colback=Green, colframe=GreenTitle,
boxrule=0.7pt, arc=0mm, title=\textbf{Formules clés}, fonttitle=\bfseries
}

\newtcolorbox{warning}{
enhanced, unbreakable, colback=Orange, colframe=OrangeTitle,
boxrule=0.7pt, arc=0mm, title=\textbf{Attention}, fonttitle=\bfseries
}

\newtcolorbox{example}{
enhanced, unbreakable, colback=Purple, colframe=PurpleTitle,
boxrule=0.7pt, arc=0mm, title=\textbf{Exemple}, fonttitle=\bfseries
}

% GOAT - brown: the idea told as simply as possible
\newtcolorbox{goat}{
enhanced, unbreakable, colback=Goat, colframe=GoatTitle,
boxrule=0.7pt, arc=0mm, title=\textbf{Expliqué à une chèvre}, fonttitle=\bfseries
}

% RECAP - violet: what to know for the exam, at the end of the week
\newtcolorbox{recap}{
enhanced, breakable, colback=Purple, colframe=RecapTitle,
boxrule=1pt, arc=0mm, title=\textbf{Récap pour l'examen}, fonttitle=\bfseries
}

% Screenshot of a slide of the course: \slide{cNN-PP}
\newcommand{\slide}[2][0.95]{%
\begin{center}
\includegraphics[width=#1\linewidth]{../images/course/#2}
\end{center}}


%=========================================================
% DOCUMENT
%=========================================================

\begin{document}

\thispagestyle{fancy}

\begin{multicols}{2}

%=========================================================
\section{Panorama des céramiques}

\begin{definition}
Familles : \textbf{naturelles} (calcaire, grès, granite, glace) ; \textbf{argiles} poreuses et
vitrifiées (porcelaine, briques) ; \textbf{verres} (sodocalcique 70 SiO$_2$/10 CaO/15 Na$_2$O,
borosilicaté) ; \textbf{ciments et bétons} ; \textbf{techniques} (Al$_2$O$_3$, SiC, Si$_3$N$_4$,
ZrO$_2$), plus fibres, revêtements, abrasifs.
\end{definition}

\begin{important}
Faites d'éléments abondants (O, Si, Al, Ca\ldots). Rigides, dures, densité moyenne, réfractaires,
inertes, souvent bon marché ; certaines recyclables (verre), d'autres non (béton). Mais
\textbf{ténacité très faible} : \textbf{fragiles}. La limite d'élasticité, énorme, ne sert à rien :
elles cassent avant.
\end{important}

%=========================================================
\section{Indices de performance}

\begin{formula}
\textbf{Barreau en traction} : $\frac{F}{\pi r^2} = E\frac{\delta}{L}$, $M = \pi r^2L\rho$
$\Rightarrow M = \frac{L^2F}{\delta}\frac{\rho}{E}$ : maximiser $\mathbf{E/\rho}$. \\
\textbf{Poutre pleine} ($t\times t$) en flexion : $\delta = \frac{4l^3F}{Et^4}$, $M = lt^2\rho$
$\Rightarrow M \propto \rho/E^{1/2}$ : maximiser $\mathbf{E/\rho^2}$. \\
\textbf{Tube mince} ($r$ fixé) en flexion : $M \propto \rho/E$ : $\mathbf{E/\rho}$.
\end{formula}

\slide[0.85]{c02-21}

\begin{important}
En flexion, rigidité $\propto t^4$, masse $\propto t^2$ : épaissir est efficace, la densité compte
\og deux fois\fg{}. Les céramiques techniques sont en haut des cartes $E$--$\rho$, mais en bas de
la carte ténacité.
\end{important}

\begin{goat}
Très rigides et légères pour leur rigidité, résistantes à la chaleur et à la corrosion, mais la
moindre fissure les casse. Pour choisir un matériau, on regarde un mélange de propriétés qui
dépend de la forme de la pièce.
\end{goat}

%=========================================================
\section{Structure des céramiques}

\begin{definition}
Liaisons \textbf{iono-covalentes} : métal + non-métal $\to$ dominante \textbf{ionique} (MgO,
Al$_2$O$_3$) ; non-métaux $\to$ dominante \textbf{covalente} (SiC, Si$_3$N$_4$) ; souvent mixte.
\end{definition}

\begin{important}
\textbf{Cristaux ioniques} : anions (gros) en empilement dense CFC/HC (ou cubique simple),
cations (petits) dans les \textbf{sites interstitiels}, choisis \textbf{un peu trop petits} pour
que les anions ne se touchent pas ; neutralité. Coordinence 4 (tétraèdre), 6 (octaèdre), 8
(cube) selon le rapport des rayons.
\end{important}

\slide[0.8]{c02-36}

Exemples : NaCl/MgO (sel gemme), CsCl, ZnS, CaF$_2$, BaTiO$_3$ (pérovskite), MgAl$_2$O$_4$
(spinelle), Al$_2$O$_3$ (O en HC, Al dans 2/3 des sites octaédriques).

\begin{important}
\textbf{Covalents} : liaisons dirigées, structures peu denses. \textbf{Silice} : tétraèdre
SiO$_4$ ; quartz, tridymite, cristobalite selon $T$ ; plus d'O/Si $\Rightarrow$ réseau 3D
$\to$ feuillets $\to$ chaînes $\to$ tétraèdres isolés. Feuillets (mica, talc, argiles) hydratables
: ils glissent, \textbf{l'argile humide est malléable}.
\end{important}

\begin{warning}
SiO$_2$--Al$_2$O$_3$, MgO--Al$_2$O$_3$ : diagrammes \textbf{quasi-binaires}, coupes de ternaires,
valables seulement à stœchiométrie d'oxygène respectée.
\end{warning}

\begin{goat}
Les gros ions s'empilent comme des oranges, les petits se glissent dans les trous. La silice est
faite de petites pyramides qu'on assemble en blocs, en feuilles ou en chaînes ; les feuilles
d'argile mouillées glissent.
\end{goat}

%=========================================================
\section{Récap de la semaine}

\begin{recap}
\begin{itemize}
\item 5 familles ; rigides, dures, réfractaires, mais fragiles.
\item Indices : $E/\rho$ (traction, tube), $E/\rho^2$ (poutre pleine en flexion) ; savoir les
      démontrer.
\item Ioniques : règles de voisinage, sites interstitiels un peu trop petits, coordinence.
\item Covalentes : dirigées ; silice et silicates (3D, feuillets, chaînes) ; argile + eau.
\item Quasi-binaires = coupes de ternaires.
\end{itemize}
\end{recap}

\end{multicols}

\end{document}
